\title{量子退火算法在工程优化中的应用}
\author{杨其臻}
\date{Oct 28, 2026}
\maketitle
工程优化问题常常呈现出多目标、高维、离散与连续变量混合以及 NP-hard 的特征，这些特性使传统方法面临严峻挑战。经典算法如模拟退火、遗传算法和混合整数规划在求解大规模实例时往往受困于计算复杂度与解的质量之间难以兼顾的矛盾。随着量子计算的兴起，量子退火算法凭借其独特的量子隧穿机制，在特定组合优化场景中展现出超越经典启发式方法的潜力。本文旨在系统梳理量子退火算法的理论基础、工程建模流程与实际部署策略，并通过典型案例分析其在航空航天、智能制造、电力系统及土木工程等领域的应用前景。\par
\chapter{量子退火算法原理}
量子退火算法源于统计物理中的伊辛模型与自旋玻璃理论，其核心思想是将优化问题编码为量子系统的哈密顿量，通过缓慢改变系统参数使量子态绝热演化至基态，从而找到全局最优解。问题哈密顿量 Hp 直接编码目标函数，而驱动哈密顿量 Hd 则负责在演化初期维持量子叠加与隧穿效应。总哈密顿量随时间线性插值，形式为 H(t)=(1-s(t))Hd+s(t)Hp，其中调度函数 s(t)从 0 单调增加至 1，控制系统从纯驱动态向问题基态的过渡。量子退火与经典模拟退火的关键差异在于前者利用量子隧穿而非热噪声跨越势垒，这一特性在高且窄的能垒场景下可显著提升逃离局部最优的概率。当前主流硬件平台 D-Wave Advantage 采用 Pegasus 拓扑，其量子比特间仅存在局部连接，实际求解时需通过链式嵌入将逻辑变量映射到物理量子比特，链强参数的选取直接影响嵌入质量与解的可靠性。\par
\chapter{工程优化问题形式化}
将工程问题转化为量子退火可处理的 QUBO 或伊辛形式，需要解决变量编码、约束处理与多目标融合三类关键技术。连续变量通常采用二进制展开或对数编码映射到二元变量，精度与量子比特数量之间存在权衡，需根据问题规模与硬件极限合理设计编码粒度。约束条件通过二次惩罚项引入目标函数，例如对等式约束施加λ(∑ xi − b)² 形式的惩罚，λ的取值需在可行域保持与目标函数量级相当。拉格朗日松弛方法则将约束乘子作为外层迭代变量，在经典求解器中动态更新，从而在量子退火调用间隙实现约束满足度的逐步提升。多目标优化可通过加权求和或ε-约束转化为单目标 QUBO，帕累托前沿的近似采样则依赖多次不同权重的退火运行结果进行后处理。\par
\chapter{典型工程场景与案例}
在航空航天领域，飞行器四维航迹优化问题可建模为包含风场约束的时空网络流问题，通过 QUBO 形式描述航段选择与燃油消耗目标；卫星星座任务调度则需同时考虑可见窗口、能量约束与数据下行优先级，量子退火在数百变量规模下展现出与商用求解器相当的解质量。智能制造中的柔性作业车间调度可编码为带有机器选择与工序排序双重决策的 QUBO 模型，实验表明混合量子-经典流程在中等规模实例上可获得优于遗传算法的 makespan。电力系统机组组合问题涉及启停状态与出力水平的混合整数决策，量子退火通过对数编码处理连续出力变量，配合经典拉格朗日松弛外层迭代，在 IEEE 标准测试系统上验证了其在旋转备用约束下的经济性优势。结构工程中的桁架拓扑优化可表述为最小重量-最大刚度双目标问题，采用二次惩罚处理应力与位移约束，量子退火生成的拓扑方案在经典有限元复核中表现出与 SIMP 方法接近的性能。\par
\chapter{混合量子-经典工作流}
受限于当前硬件的量子比特数量与连通性，大规模工程问题通常采用混合量子-经典工作流进行求解。预处理阶段通过变量固定、聚类或对称约简降低有效问题规模；分解策略则将原问题划分为多个子 QUBO，由 D-Wave Hybrid Solver 或开源 qbsolv 框架交替调用量子退火与经典局部搜索。嵌入过程中链强参数需根据问题能量尺度动态调节，过低导致链断裂，过高则压缩问题哈密顿量有效间隙。误差缓解技术包括自旋反转后处理与多数投票，统计上可降低读出噪声对基态判断的影响。云端部署示例中，AWS Bracket 或 D-Wave Leap 云服务提供 REST 接口，边缘节点负责问题建模与结果后处理，形成低延迟的工程决策闭环。\par
\chapter{性能评估与 benchmarking}
评估量子退火在工程优化中的实际价值需建立时间-质量权衡曲线，常用指标包括达到目标精度的中位退火时间 TTS 与总计算时间 TTT。通过与 Gurobi、CPLEX 在相同 QUBO 实例上的对比实验，可量化量子退火在不同规模与约束密度下的相对优势。统计显著性检验采用 Wilcoxon 秩和检验或 Bootstrap 重采样，确保结论不受单次运行随机性影响。可重复性要求完整记录嵌入参数、退火时间表、随机种子及后处理流程，以便第三方在相同或不同硬件平台上复现结果。\par
\chapter{挑战与未来方向}
硬件层面，连通性不足与噪声退相干仍是制约量子退火大规模应用的核心瓶颈，下一代量子退火机需在千比特以上规模与全连通拓扑上取得突破。算法层面，近似比的理论保证尚不完善，超参数对解质量的敏感性要求工程用户具备深厚的调参经验。生态建设方面，标准化 QUBO 建模语言与容器化接口的缺失导致不同平台间迁移成本高昂。变分退火 VAQAO 与量子近似优化算法 QAOA 的融合研究正在兴起，目标是结合变分原理与绝热思想，在浅层量子电路与退火硬件间寻找折中方案。量子启发经典算法如数字模拟退火与基于张量网络的近似方法，也为在经典硬件上模拟量子优势提供了新思路。\par
\chapter{结论}
量子退火算法在工程优化中的定位是经典方法的有效补充，而非全面替代。在短期内，其商业落地价值集中于数百至数千变量的组合优化子问题，通过混合工作流可显著缩短设计迭代周期。工程教育需加强量子信息与优化理论的跨学科融合，培养既懂领域物理又掌握量子建模的复合型人才，方能真正释放量子退火在复杂工程系统中的应用潜力。\par
\chapter{参考文献与资源}
Kadowaki 与 Nishimori 于 1998 年提出的量子退火框架奠定了该领域理论基础，Farhi 等人关于绝热量子计算的系列工作则为算法分析提供了严格数学工具。D-Wave 公司发布的白皮书详细阐述了 Pegasus 与 Zephyr 拓扑的工程实现细节。开源工具链方面，Ocean SDK 提供与 D-Wave 硬件的原生接口，PyQUBO 负责 QUBO 建模与符号推导，Dimod 则实现统一采样器抽象。在线 benchmark 数据集与竞赛平台如 Qiskit Application Repository 与 Quantum Optimization Challenge，为研究者提供了标准测试实例与性能排行榜。\par
